\hypertarget{budget}{%
\section{5. Budget}\label{budget}}

This chapter presents the budget of the Namibia 2026 deployment as the
reference instance of the recipe described in Chapters 3 and 4. The
budget is structured in five blocks --- network hardware, endpoint
hardware, software, engineering effort, and travel and logistics ---
followed by a discussion of funding sources, a per-beneficiary
breakdown, and a comparison with the cost of the equivalent
new-equipment, commercial-deployment alternative.

Two notes apply to every figure that follows. First, all prices are
expressed in \textbf{euros (€) including VAT} unless stated otherwise.
Second, several entries are marked \textbf{(TBD: confirm against
receipts)}; these are the items for which the field log contains the
canonical figure but which had not been consolidated into the project
ledger at the time of submission. The estimates given are based on the
public list price at the time of purchase (October 2025 to February
2026) and the figures recorded in the AUCOOP project tracker; they are
within \textasciitilde10 \% of the receipted values and the ranking of
cost categories is robust to that uncertainty.

\begin{center}\rule{0.5\linewidth}{0.5pt}\end{center}

\hypertarget{network-hardware}{%
\subsection{5.1 Network hardware}\label{network-hardware}}

\begin{longtable}[]{@{}lrrr@{}}
\toprule\noalign{}
Item & Unit cost (€) & Qty & Subtotal (€) \\
\midrule\noalign{}
\endhead
\bottomrule\noalign{}
\endlastfoot
Cudy WR3000E mesh-capable AP & \textasciitilde55 & 7 &
\textasciitilde385 \\
NanoPi R6S gateway (8 GB) & \textasciitilde110 & 1 &
\textasciitilde110 \\
TP-Link 8-port Gigabit PoE switch & \textasciitilde95 & 1 &
\textasciitilde95 \\
Cat-6 patch cables and connectors & \textasciitilde3 & 30 &
\textasciitilde90 \\
Outdoor-rated Ethernet (per metre) & \textasciitilde1.20 & 50 &
\textasciitilde60 \\
24 V / 12 V PoE injectors and splitters & \textasciitilde12 & 4 &
\textasciitilde48 \\
Lithium UPS (small, for gateway + switch) & \textasciitilde120 & 1 &
\textasciitilde120 \\
Wall-mounting hardware, conduit, ties & (TBD) & --- &
\textasciitilde50 \\
\textbf{Subtotal --- Network hardware} & & &
\textbf{\textasciitilde958} \\
\end{longtable}

The figures above describe what was procured for the deployment; they do
not include the spare AP and the spare PoE injector kept on the bench in
Barcelona, both of which were existing AUCOOP stock. The Cudy WR3000E is
the single most numerous item and is therefore the budget line most
sensitive to a change of recipe; §3.A.4 documents the rationale for that
choice and §6.2 the environmental side of it.

\hypertarget{endpoint-hardware}{%
\subsection{5.2 Endpoint hardware}\label{endpoint-hardware}}

\begin{longtable}[]{@{}
  >{\raggedright\arraybackslash}p{(\columnwidth - 6\tabcolsep) * \real{0.2000}}
  >{\raggedleft\arraybackslash}p{(\columnwidth - 6\tabcolsep) * \real{0.2667}}
  >{\raggedleft\arraybackslash}p{(\columnwidth - 6\tabcolsep) * \real{0.2667}}
  >{\raggedleft\arraybackslash}p{(\columnwidth - 6\tabcolsep) * \real{0.2667}}@{}}
\toprule\noalign{}
\begin{minipage}[b]{\linewidth}\raggedright
Item
\end{minipage} & \begin{minipage}[b]{\linewidth}\raggedleft
Unit cost (€)
\end{minipage} & \begin{minipage}[b]{\linewidth}\raggedleft
Qty
\end{minipage} & \begin{minipage}[b]{\linewidth}\raggedleft
Subtotal (€)
\end{minipage} \\
\midrule\noalign{}
\endhead
\bottomrule\noalign{}
\endlastfoot
Refurbished Lenovo ThinkPad T460 / X260 & 0 (donated) & 9 & 0 \\
Replacement RAM (8 GB DDR4 modules) & \textasciitilde25 & 4 &
\textasciitilde100 \\
Replacement SSDs (240 GB) & \textasciitilde30 & 3 & \textasciitilde90 \\
Replacement batteries (genuine / refurbished) & \textasciitilde45 & 2 &
\textasciitilde90 \\
Replacement palmrests / keyboards & (TBD) & 1--2 & \textasciitilde70 \\
External USB-A sticks (16 GB, for AUCOOP-image and PXE seed) &
\textasciitilde6 & 5 & \textasciitilde30 \\
External USB SSD (500 GB, image staging) & \textasciitilde55 & 1 &
\textasciitilde55 \\
Cleaning, thermal paste, screws & (TBD) & --- & \textasciitilde25 \\
\textbf{Subtotal --- Endpoint hardware} & & &
\textbf{\textasciitilde460} \\
\end{longtable}

The headline is the zero in the first row: the nine endpoints themselves
were donated through Labdoo and NexTReT and entered the project at no
cash cost. The €460 of cash spend on the endpoint side is therefore
entirely \emph{refurbishment} spend --- the parts and consumables
required to take the donations from incoming-equipment status to
deployed at school. The ratio of refurbishment cost to donated value is
the operationally relevant figure: at a notional secondary-market price
of €150 per refurbished T460, the nine machines represent roughly €1 350
of asset value mobilised at €460 of cash, a leverage of
\textasciitilde3:1.

\hypertarget{software}{%
\subsection{5.3 Software}\label{software}}

All software used in the deployment is free and open source: OpenWrt,
Linux Mint Cinnamon, ISC \texttt{dhcpd}, \texttt{tftpd-hpa}, NFS, GRUB,
\texttt{pxelinux}, Clonezilla, \texttt{partclone}, OnlyOffice (community
edition), Zensical, MkDocs Material, GitHub Actions on a public
repository.

\textbf{Subtotal --- Software: €0.}

The zero is itself part of the recipe. Section 3.B.5 (AUCOOP image),
§3.A.4 (OpenWrt) and §3.C (handbook tooling) all rely on software whose
licensing imposes no per-unit cost and whose source is auditable. The
economic implication is that the recipe's cash cost scales with hardware
quantity only; the software cost of deploying a second school is the
same as the first, namely zero.

\hypertarget{engineering-effort}{%
\subsection{5.4 Engineering effort}\label{engineering-effort}}

The thesis budget accounts for engineering effort using indicative
hourly rates rather than market rates; the real cost is borne by the
university (advisor time) and by the student (own time). The table is
included for completeness and to make the implicit investment legible.

\begin{longtable}[]{@{}lrrr@{}}
\toprule\noalign{}
Role & Hours & Rate (€/h) & Subtotal (€) \\
\midrule\noalign{}
\endhead
\bottomrule\noalign{}
\endlastfoot
Student engineer (M. Jover, this thesis) & \textasciitilde720 & 25 &
\textasciitilde18 000 \\
Student engineer (J. Motje, companion thesis) & \textasciitilde720 & 25
& \textasciitilde18 000 \\
Advisor (E. Vidal) & \textasciitilde60 & 60 & \textasciitilde3 600 \\
Co-advisor (S. Giménez) & \textasciitilde30 & 40 & \textasciitilde1
200 \\
AUCOOP volunteer hours (intake, logistics) & \textasciitilde80 & 20 &
\textasciitilde1 600 \\
\textbf{Subtotal --- Engineering effort} & & & \textbf{\textasciitilde42
400} \\
\end{longtable}

Two observations on this table. First, the engineering line is by a very
large margin the dominant cost block, dwarfing the network and endpoint
hardware combined; this is consistent with the framing of §2.1, which
characterises community networks as substituting capital for human
expertise. Second, the engineering effort is split across two theses;
the figure double-counts the advisor hours that were spent on shared
meetings. A pessimistic single-thesis attribution would be roughly half
of the figure above plus the full advisor share.

\hypertarget{travel-and-logistics-namibia-field-deployment}{%
\subsection{5.5 Travel and logistics (Namibia field
deployment)}\label{travel-and-logistics-namibia-field-deployment}}

\begin{longtable}[]{@{}lr@{}}
\toprule\noalign{}
Item & Cost (€) \\
\midrule\noalign{}
\endhead
\bottomrule\noalign{}
\endlastfoot
Return flight Barcelona ↔ Windhoek (×2 students) & \textasciitilde2
400 \\
Domestic transport Windhoek ↔ Gochas (4×4 hire, fuel) &
\textasciitilde600 \\
Lodging at Gochas (10 nights, ×2) & \textasciitilde500 \\
Per-diem (food, water, local consumables) & \textasciitilde400 \\
Travel insurance (×2) & \textasciitilde120 \\
Equipment shipping (Foundawtion container share) & \textasciitilde200 \\
Vaccinations and visa & \textasciitilde180 \\
\textbf{Subtotal --- Travel and logistics} & \textbf{\textasciitilde4
400} \\
\end{longtable}

The flights alone account for more than half of this block and are also
the single largest contributor to the project's transport carbon
footprint (§6.7). The container share for equipment shipping is reduced
because Foundawtion runs the container as a recurring operation and the
project pays only the marginal cost of the volume occupied.

\hypertarget{funding-sources}{%
\subsection{5.6 Funding sources}\label{funding-sources}}

The project's cash needs are met from three sources, none of which flow
through the student authors directly.

\begin{itemize}
\tightlist
\item
  \textbf{UPC CCD (Centre de Cooperació per al Desenvolupament)}:
  \textbf{5 600 €} competitive grant covering travel, network hardware
  procurement, and refurbishment consumables. The grant is the largest
  single source and is the reason the recipe is reproducible inside UPC:
  subsequent cohorts can apply for the same grant under the same call.
\item
  \textbf{Labdoo}: nine laptops in kind (T460 and X260 generation),
  delivered already provenanced through the Labdoo tracking system. No
  cash flow, \textasciitilde€1 350 of asset value.
\item
  \textbf{NexTReT}: three laptops in kind plus two mini-PC servers used
  for the services thesis; \textasciitilde€500 of asset value
  attributable to the hardware-side scope.
\item
  \textbf{AUCOOP}: coordination, intake-and-staging space at the UPC
  Campus Nord, and the volunteer hours costed in §5.4. No cash flow.
\end{itemize}

The grant covers approximately 100 \% of the cash spend recorded in
§§5.1, 5.2 and 5.5 (\textasciitilde6 000 € against the 5 600 € grant,
with a \textasciitilde400 € overshoot absorbed by AUCOOP's own funds).
The engineering effort of §5.4 is explicitly not funded; it is
contributed in exchange for the academic credit of the two theses.

\hypertarget{total}{%
\subsection{5.7 Total}\label{total}}

\begin{longtable}[]{@{}lr@{}}
\toprule\noalign{}
Block & € \\
\midrule\noalign{}
\endhead
\bottomrule\noalign{}
\endlastfoot
Network hardware (§5.1) & \textasciitilde958 \\
Endpoint hardware (§5.2) & \textasciitilde460 \\
Software (§5.3) & 0 \\
Engineering effort (§5.4) & \textasciitilde42 400 \\
Travel and logistics (§5.5) & \textasciitilde4 400 \\
\textbf{TOTAL --- full project} & \textbf{\textasciitilde48 218} \\
\textbf{TOTAL --- cash only (excludes engineering)} &
\textbf{\textasciitilde5 818} \\
\end{longtable}

The two totals tell different stories. The full-project figure of
\textasciitilde48 k€ reflects the true societal cost of the deployment
if every hour of student and advisor time were paid at indicative rates;
this is the figure that makes the project comparable to a commercial
deployment. The cash-only figure of \textasciitilde5.8 k€ is what an
organisation with access to a similar volunteer pool actually has to
raise; this is the figure that makes the recipe replicable inside AUCOOP
and similar associations, and it is the one against which §5.9 compares
the commercial alternative.

\hypertarget{cost-per-beneficiary}{%
\subsection{5.8 Cost per beneficiary}\label{cost-per-beneficiary}}

The N Mutschuana Primary School serves approximately \textbf{300
students}. Distributed evenly:

\begin{itemize}
\tightlist
\item
  \textbf{Full-project cost per student}: \textasciitilde48 218 ÷ 300 ≈
  \textbf{161 € / student}.
\item
  \textbf{Cash cost per student}: \textasciitilde5 818 ÷ 300 ≈
  \textbf{19 € / student}.
\item
  \textbf{Cash cost per student per year of service} (assuming a
  conservative 5-year service life): \textasciitilde{}\textbf{4 € /
  student / year}.
\end{itemize}

The €4-per-student-per-year figure is the one that lands. It includes a
connected wireless network covering the school and a fleet of nine
endpoints, and it is well within the range of cost-effectiveness
benchmarks for ICT-in-education programmes reported in the ICT4D
literature surveyed in §2.3.

\hypertarget{comparison-with-the-commercial-equivalent-alternative}{%
\subsection{5.9 Comparison with the commercial-equivalent
alternative}\label{comparison-with-the-commercial-equivalent-alternative}}

The point of the budget is partly comparative. A commercial alternative
--- new business-grade laptops, vendor mesh access points, contracted
installer, three-year support --- would price the same deployment as
follows.

\begin{longtable}[]{@{}
  >{\raggedright\arraybackslash}p{(\columnwidth - 4\tabcolsep) * \real{0.2727}}
  >{\raggedleft\arraybackslash}p{(\columnwidth - 4\tabcolsep) * \real{0.3636}}
  >{\raggedleft\arraybackslash}p{(\columnwidth - 4\tabcolsep) * \real{0.3636}}@{}}
\toprule\noalign{}
\begin{minipage}[b]{\linewidth}\raggedright
Block
\end{minipage} & \begin{minipage}[b]{\linewidth}\raggedleft
Refurbished + open recipe
\end{minipage} & \begin{minipage}[b]{\linewidth}\raggedleft
Commercial-equivalent
\end{minipage} \\
\midrule\noalign{}
\endhead
\bottomrule\noalign{}
\endlastfoot
Network hardware & \textasciitilde960 € & \textasciitilde3 500 € (Aruba
Instant On / equivalent) \\
Endpoints & \textasciitilde460 € (parts) & \textasciitilde7 200 € (9 ×
\textasciitilde800 € new business laptops) \\
Software / licensing & 0 € & \textasciitilde1 000 € (Windows +
productivity, 9 seats) \\
Installation labour & (volunteer) & \textasciitilde2 500 € (contractor,
5 days × 2 people) \\
3-year support contract & (community) & \textasciitilde1 800 € \\
\textbf{Cash subtotal} & \textbf{\textasciitilde1 420 €} &
\textbf{\textasciitilde16 000 €} \\
Travel and logistics & \textasciitilde4 400 € & \textasciitilde4 400 €
(unchanged) \\
\textbf{Total cash} & \textbf{\textasciitilde5 820 €} &
\textbf{\textasciitilde20 400 €} \\
\end{longtable}

The commercial alternative is approximately \textbf{3.5× more expensive
in cash terms} and carries an additional environmental cost in
manufacturing-phase emissions (quantified in §6.2). The community-recipe
approach trades that cash and that footprint for volunteer engineering
effort and for the operational dependency on the handbook of §3.C --- a
dependency that the validation chapter (§4.4) shows the handbook is now
mature enough to carry.

The comparison should not be over-read. A commercial deployment buys
things the community recipe does not: a single-throat-to-choke support
contract, an SLA, vendor liability, and faster mean-time-to- repair on
hardware failures. Whether those properties are worth the 3.5× premium
depends on the deploying organisation's risk profile; for AUCOOP and its
partner sites, the trade-off has historically favoured the recipe
described in this thesis, and this budget makes the magnitude of that
trade-off explicit.
